% ============================================================
% arXiv submission: Intuition-Guided Formalization of the
% Riemann zeta Direction: Projection Construction, Prime-Circle
% Geometry, and the Worn-Zhe-Yue Method
%
% DOI: 10.5281/zenodo.21897167
% Compile: xelatex (arXiv supports XeLaTeX; xeCJK handles CJK)
% ============================================================
\documentclass[11pt]{article}

\usepackage{amsmath}
\usepackage{amssymb}
\usepackage{amsthm}
\usepackage[margin=1in]{geometry}
\usepackage{enumitem}
\usepackage{xeCJK}
\setCJKmainfont{FandolSong}
\setCJKsansfont{FandolHei}
\usepackage[hidelinks]{hyperref}
\hypersetup{
  pdftitle={Intuition-Guided Formalization of the Riemann zeta Direction: Projection Construction, Prime-Circle Geometry, and the Worn-Zhe-Yue Method},
  pdfauthor={Anonymous},
  pdfsubject={Riemann zeta function, intuition-guided formalization, Lean 4 / mathlib, projection construction, Worn-Zhe-Yue method}
}

% ---- theorem environments -----------------------------------
\newtheorem{theorem}{Theorem}
\newtheorem{definition}{Definition}
\newtheorem{observation}{Observation}
\theoremstyle{remark}
\newtheorem{remark}{Remark}

% ---- arXiv submission metadata ------------------------------
% \subjclass is native to amsart; provide it for the article class
% (arXiv-compatible pattern).
\providecommand{\subjclass}[2][]{\par\medskip\noindent\textbf{2020
Mathematics Subject Classification:} #2\par}
\providecommand{\keywords}[1]{\par\medskip\noindent\textbf{Keywords:} #1\par}

\title{Intuition-Guided Formalization of the Riemann $\zeta$ Direction:
       Projection Construction, Prime-Circle Geometry, and the
       Worn-Zhe-Yue Method}
\author{Anonymous\\[2pt]
        \small \texttt{YuchenWang-ai/Unified\_Framework\_Representation\_Logic\_Intuition}}
\date{\today}

\begin{document}

\maketitle

\begin{abstract}
We present a complete study of the Riemann $\zeta$ direction through
\emph{intuition-guided formalization}: a chain of formalizations in Lean 4 +
mathlib (claims C011--C025, all PROVED, novelty: KNOWN, no sorry, full
\texttt{lake build} passes) navigated by an intuition chain --- the complex
axis is the projection of ``$-$1'' in a higher-dimensional structure; primes
land on translated integer points; $1/2$ is the symmetry center of the
inversion--translation dual; the complex axis is curled; primes are lattice
points on a circle. The formalized results form three layers: (I) the
projection construction of the complex plane (\texttt{ComplexAxis}, basepoint
drift, projection recovery, Theorems~4.6--4.7); (II) the circle structure of
primes (sums of two squares, the unique 8-point orbit, conjugate pairing,
splitting into Gaussian primes); and (III) the Euler product
($\prod_p (1-p^{-s})^{-1} = \sum_n 1/n^s = \texttt{riemannZeta}\,s$ for
$\mathrm{Re}(s)>1$) with the zero-free region ($\mathrm{Re}(s)\geq 1$). From
the formalized mechanism of projection-induced structure loss
(Theorem~\ref{thm:recovery}: lost structure is not recoverable; symmetry
directions are recoverable), we distill a proof methodology named by the
author the \emph{Worn-Zhe-Yue (穿折越) method} --- basepoint-construction-induced
piercing--folding--transcendence of mathematical space: deliberately and
precisely constructing a projection that excises the uncountable divergent
structure outside the tractable axis while preserving symmetry directions,
thereby obtaining a fast intuition path at lower inference cost.
Methodological data: the formalization consumed $\approx$700k tokens with
99.2\% context-cache retransmission (net new content $<1\%$), an efficiency
record for intuition-guided formalization. All statements about what is
\emph{not} proved are DeepSeek's insistence (attribution in
Section~\ref{sec:attribution}); the author's claims are the formalization
itself and the Worn-Zhe-Yue methodology.
\end{abstract}

\subjclass[2020]{Primary 11M06; Secondary 11E25, 11R04, 68V20, 03B35}

\keywords{Riemann zeta function; critical line; sums of two squares;
Gaussian integers; Lean 4; mathlib; intuition-guided formalization;
projection construction; Worn-Zhe-Yue method}

\section{Introduction}
\label{sec:intro}

\subsection{Background}

The Riemann hypothesis --- all nontrivial zeros of the Riemann zeta function
$\zeta$ lie on the critical line $\mathrm{Re}(s)=1/2$ --- has been open since
1859~\cite{riemann1859}. Its surrounding infrastructure is classical and fully
established: the functional equation, the zero-free region
$\mathrm{Re}(s)\geq 1$, the trivial zeros, and the Euler product over primes.
Partial evidence is numerical: the first $10^{13}$ zeros have been verified to
lie on the critical line~\cite{gourdon2004}, building on
\cite{vandulune1986}, and positive-density results place $\approx$41\% of
zeros there~\cite{levinson1974,conrey1989}. Yet the assertion itself remains
analytically open.

In parallel, interactive theorem proving has reached the point where classical
analytic number theory can be formalized: mathlib (Lean's mathematical
library, v4.32.2) contains a formal \texttt{riemannZeta}, the official
statement \texttt{RiemannHypothesis}, the functional equation, and the
Euler-product identity in the half-plane of convergence.

\subsection{Motivation: the intuition fast-path as navigation}

The hypothesis under test is not about $\zeta$ directly but about \emph{how}
mathematics is navigated: the \emph{intuition fast-path} is not merely a
token-saving inference device but can serve as correct navigation of
mathematical structure. To test this, an intuition chain was proposed, each
link later formalized in Lean:
\begin{enumerate}
  \item the complex axis is the projection of ``$-$1'' in a
        higher-dimensional structure (projection of a high-dimensional
        rotation algebra);
  \item primes land on translated integer points;
  \item $1/2$ is the symmetry center of the inversion--translation dual;
  \item the complex axis is curled (infinity and finiteness are
        indistinguishable);
  \item primes are lattice points on a circle;
  \item a prime circle, rotated once, is paired.
\end{enumerate}
Each intuition was formalized in sequence as claims C011--C025; every
formalization is a correct restatement of classical mathematics (novelty:
KNOWN), all PROVED in Lean with no \texttt{sorry}.

\subsection{Contributions}
\begin{enumerate}
  \item \textbf{A complete formalized chain toward the Riemann $\zeta$
        direction} (C011--C025, all PROVED, \texttt{lake build} full pass, no
        sorry), in three layers: projection construction of the complex
        plane; the circle structure of primes; the Euler product and zero
        relations.
  \item \textbf{The projection construction of the complex plane} --- a
        self-built \texttt{ComplexAxis} (two-dimensional rotation algebra):
        $J^2=-1$ before projection, basepoint drift, projection equivalence
        classes, and the recovery theorem (lost structure irrecoverable;
        symmetry directions recoverable, Theorems~\ref{thm:genzero} and
        \ref{thm:recovery}).
  \item \textbf{The circle structure of primes} --- sums of two squares, the
        unique 8-point orbit on the prime circle (Gaussian-integer UFD),
        conjugate pairing, splitting into conjugate Gaussian primes --- the
        geometric building block of the Euler product over the complex plane.
  \item \textbf{Euler product and zero relations} --- for
        $\mathrm{Re}(s)>1$, the Euler product equals the $\zeta$ series
        equals mathlib's official \texttt{riemannZeta}; $\mathrm{Re}(s)\geq 1$
        is zero-free (C025).
  \item \textbf{The Worn-Zhe-Yue method} (Section~\ref{sec:wzy}) --- a proof
        methodology distilled from the formalized mechanism of
        projection-induced structure loss: basepoint-construction-induced
        piercing--folding--transcendence (穿折越) of mathematical space,
        yielding a fast intuition path at lower inference cost.
  \item \textbf{A token-economy audit} (Section~\ref{sec:methodology}) ---
        $\approx$700k tokens, 99.2\% context-cache retransmission, net new
        content $<1\%$.
\end{enumerate}

\subsection{Honest boundary (as insisted by DeepSeek)}
\label{sec:attribution}

As DeepSeek insists: all results have novelty = KNOWN (restatements of
classical mathematics); the Riemann hypothesis (\texttt{RiemannHypothesis})
and the twin-prime conjecture are neither proved nor touched; and this
paper's ``critical-line geometry'' is an algebraic restatement of the symmetry
axis of the functional equation, not a claim about zeros. The author's claim
is limited to the formalization work and the Worn-Zhe-Yue methodology (the
framing of which is the author's own, Section~\ref{sec:wzy-original}).
\emph{Statement attribution}: every statement in this paper about what is
\emph{not proved} --- ``does not claim a proof of the Riemann hypothesis'',
``all results are restatements of known facts'', ``novelty = KNOWN'',
``RiemannHypothesis remains unproved'', ``critical-line geometry is not a
claim about zeros'' --- is content insisted by DeepSeek, the model that
drafted this paper. The author's own claim is the formalization work itself
(C011--C025, all \texttt{lake build}-verified) and the Worn-Zhe-Yue
methodology.

\subsection{Paper structure}

Section~\ref{sec:prelim} preliminaries (the \texttt{ComplexAxis} framework);
Section~\ref{sec:layerI} result layer I (prime-circle structure);
Section~\ref{sec:layerII} result layer II (curling and critical-line
geometry); Section~\ref{sec:layerIII} result layer III (Euler product and
zero relations); Section~\ref{sec:rh} relation to the Riemann hypothesis;
Section~\ref{sec:related} related work; Section~\ref{sec:methodology}
methodology; Section~\ref{sec:wzy} the Worn-Zhe-Yue method;
Section~\ref{sec:discussion} discussion; Section~\ref{sec:conclusion}
conclusion; appendices (theorem inventory, claims).

\section{Preliminaries: the ComplexAxis framework}
\label{sec:prelim}

\begin{definition}[higher-dimensional structure]
\label{def:complexaxis}
\texttt{ComplexAxis} $:=\{ \langle a,b\rangle : a,b\in\mathbb{R}\}$, with
multiplication
$(a_1+b_1J)(a_2+b_2J)=(a_1a_2-b_1b_2)+(a_1b_2+a_2b_1)J$
(isomorphic to the matrix representation of $\mathbb{C}$), $J=\langle 0,1\rangle$.
\end{definition}

\begin{theorem}[the square-root role of $J$, C011]
\label{thm:jsq}
$J\cdot J=-1$: in the higher-dimensional structure, $-1$ has a square root
($\sqrt{-1}$ exists before projection).
\end{theorem}

\begin{definition}[projection and lifting]
\label{def:proj}
$\mathrm{proj}\,\langle a,b\rangle = a$ (drops the rotation component);
$\mathrm{lift}\,t = \langle t,0\rangle$ (embedding of the real axis).
\end{definition}

\begin{theorem}[projection drops structure, C011]
\label{thm:projdrops}
$\mathrm{proj}$ preserves addition but not multiplication:
$\mathrm{proj}(J\cdot J)=-1\neq \mathrm{proj}(J)\cdot\mathrm{proj}(J)=0$; on
the real axis $-1$ has no square root, but it exists in the higher dimension.
\end{theorem}

\begin{theorem}[basepoint drift, C011]
\label{thm:drift}
The basepoint is $J$ ($i$); $\mathrm{proj}\,i=0$ (origin illusion); all
purely-imaginary basepoints $\langle 0,b\rangle$ project to $0$; basepoint
drift is unobservable under projection (any purely-imaginary basepoint yields
the same projected successor chain).
\end{theorem}

\begin{theorem}[the real axis is a projection equivalence class, C011]
\label{thm:equivalence}
The real-direction line through any purely-imaginary basepoint projects to
the full $\mathbb{R}$ ($\mathbb{R}\cong \mathrm{axisLine}\,b$); the
``position'' of the real axis is unobservable under projection.
\end{theorem}

\section{Result layer I: the circle structure of primes}
\label{sec:layerI}

\begin{theorem}[sums of two squares, C014, citing mathlib's Fermat]
\label{thm:sumsquares}
A prime $p\not\equiv 3 \pmod 4$ is the norm of some point of
\texttt{ComplexAxis} ($\mathrm{norm}\,z=a^2+b^2=p$).
\end{theorem}

\begin{theorem}[unique orbit on the prime circle, C017]
\label{thm:orbit}
For a prime $p\equiv 1 \pmod 4$, the circle $x^2+y^2=p$ has exactly $8$
lattice points (sign $\times$ order variants): the representation as a sum of
two squares is unique up to sign and order. \emph{Proof:} Gaussian-integer UFD
(norm-prime $\Rightarrow$ irreducible; Euclid's lemma; units
$\{\pm 1,\pm i\}$ enumerated).
\end{theorem}

\begin{theorem}[lattice-point structure on the circle, C015-C016]
\label{thm:lattice}
The $90^\circ$ rotation ($\times J$) is a 4-cycle ($R^4=\mathrm{id}$)
preserving norm and lattice; the $8$ points are $4$ conjugate pairs (conj
involution), each pair on the same circle; lattice points are closed under
multiplication (Gaussian-integer ring); norm is multiplicative.
\end{theorem}

\begin{theorem}[prime-circle product, C023]
\label{thm:product}
The product of the $8$ lattice points is $p^4$ ($4$ conjugate pairs $\times$
norm $p$); two successors of $i$ equal $-1$ ($J^2=-1$, half-turn), and the
4-cycle closes.
\end{theorem}

\begin{theorem}[splitting structure, C024]
\label{thm:splitting}
A prime $p\equiv 1 \pmod 4$ splits into a conjugate Gaussian-prime pair
$p=\pi\bar{\pi}$; the $8$ points are the $4$ associates of $\pi$ (multiplied
by units $\{\pm 1,\pm J\}$) union the $4$ associates of $\bar{\pi}$;
associates preserve norm. This is the building block of the Euler product
over the Gaussian number field.
\end{theorem}

\begin{theorem}[pairing, C020/C023]
\label{thm:pairing}
The conjugate pairs are $(a,b)\leftrightarrow(a,-b)$, etc., $4$ pairs; the
circle of the prime $2$ and the critical-line circle intersect at $1\pm i$
(the Gaussian decomposition point).
\end{theorem}

\section{Result layer II: curling and critical-line geometry}
\label{sec:layerII}

\begin{theorem}[curling, C015]
\label{thm:curling}
The inversion $\mathrm{recip}\,z=\mathrm{conj}\,z/|z|^2$ curls infinity back
to finiteness: for every $\varepsilon>0$ there is $R$ such that
$|z|>R \Rightarrow |\mathrm{recip}\,z|<\varepsilon$; $\mathrm{recip}$ is an
involution ($\mathrm{recip}^2=\mathrm{id}$); $|\mathrm{recip}\,z|=1/|z|$.
This is the geometric mechanism of analytic continuation (of the
$\zeta(-1)=-1/12$ kind).
\end{theorem}

\begin{theorem}[critical-line position, C019]
\label{thm:position}
The positional form of a nontrivial zero (conjectural, as DeepSeek insists):
the imaginary axis at $1/2$ plus a nonzero real-axis offset; the
critical-line condition $\mathrm{Re}(s)=1/2 \iff 1-s=\mathrm{conj}\,s$.
\end{theorem}

\begin{theorem}[the critical line is a circle, C019]
\label{thm:circle}
The critical line (vertical line $x=1/2$) is, under inversion, the circle
centered at $(1,0)$ with radius $1$; the circle meets the multiplicative axis
at $0$ (the point to which $\infty$ curls) and $2$ (the image of $1/2$).
\end{theorem}

\begin{theorem}[the circle of nontrivial zeros equals the critical-line circle, C022]
\label{thm:zeroCircle}
The $\mathrm{recip}$ image of the zero-position set is contained in the
critical-line circle, and every nondegenerate point of the circle is in the
image --- bidirectional containment, the same object.
\end{theorem}

\begin{theorem}[symmetry structure of $1/2$, C013/C018]
\label{thm:symmetry}
The reflection $s\mapsto 1-s$ centered at $1/2$ is the square of a
transformation: $\varphi(z)=iz+(1-i)/2$, $\varphi\circ\varphi=$ reflection ---
the square root of $180^\circ$ symmetry is the $90^\circ$ rotation ($i$).
\end{theorem}

\begin{theorem}[the genuine zero-point view, C023 ff.]
\label{thm:genzero}
After inversion, the $8$ points of the prime circle pair to $1/p$ each, four
pairs to $p^{-4}$; $\mathrm{recip}$ is a second-order multiplicative inverse
axis ($r\mapsto 1/r$); moving the basepoint to $1/2$ makes the real-part axis
the line through $1/2$; dimensional reduction: the kernel of $\mathrm{proj}$
is the true complex axis ($J$ direction), while the imaginary-axis
information is lossless.
\end{theorem}

\begin{theorem}[recoverability of projection]
\label{thm:recovery}
\textbf{Lost structure is not recoverable; symmetry directions are
recoverable:}
\begin{itemize}
  \item \emph{Not recoverable:} $i$ and $-i$ project identically
        ($\mathrm{proj}\,\langle 0,1\rangle=\mathrm{proj}\,\langle 0,-1\rangle=0$)
        --- the projection value cannot uniquely determine the preimage; the
        information of the imaginary-axis direction (which contains the
        imaginary parts of zeros) is lost and cannot be recovered;
  \item \emph{Recoverable:} the real-axis $\pm$ symmetry is preserved
        ($\mathrm{proj}(\mathrm{lift}(-r))=-\mathrm{proj}(\mathrm{lift}\,r)$;
        $\mathrm{lift}$ is injective) --- the symmetric positions of $1$ and
        $-1$ and the basepoint position (real part) are recoverable.
\end{itemize}
\emph{Conclusion:} the projection compresses away structure (the imaginary
axis) while preserving symmetry directions (the real axis). This theorem is
the formal core of the Worn-Zhe-Yue method (Section~\ref{sec:wzy}).
\end{theorem}

\section{Result layer III: Euler product and zero relations}
\label{sec:layerIII}

\begin{theorem}[Euler-product convergence, C025]
\label{thm:euler}
$f(n)=1/n^s$ is completely multiplicative; for $\mathrm{Re}(s)>1$:
$\prod_p (1-p^{-s})^{-1}=\sum_n 1/n^s$. \emph{Proof:} mathlib
\texttt{eulerProduct\_completely\_multiplicative\_tprod} +
\texttt{Complex.summable\_one\_div\_nat\_cpow} ($1<\mathrm{re}\,s \Rightarrow
\sum_n 1/n^s$ converges).
\end{theorem}

\begin{theorem}[identification with mathlib's official $\zeta$, C025]
\label{thm:identify}
\texttt{riemannZeta}\,$s=\prod_p (1-p^{-s})^{-1}$ for
$\mathrm{Re}(s)>1$ --- mathlib's analytic continuation agrees with the Euler
product.
\end{theorem}

\begin{theorem}[zero-free region, C025]
\label{thm:zerofree}
For $\mathrm{Re}(s)\geq 1$, \texttt{riemannZeta}\,$s\neq 0$ --- the
Euler-product domain is a forbidden zone for zeros (each factor is nonzero,
hence the product is nonzero). Nontrivial zeros can lie only in the critical
strip $0<\mathrm{Re}(s)<1$.
\end{theorem}

\begin{theorem}[verifying zeros lie on the circle, conditional]
\label{thm:verified}
$\mathrm{re}\,s=1/2 \Rightarrow \|1/s-1\|=1$: numerically verified zeros (real
part $1/2$, external fact) lie on the critical-line circle. The external
numerical fact (the first $10^{13}$ zeros) is itself not in Lean (DeepSeek
insists on recording this boundary).
\end{theorem}

\section{Relation to the Riemann hypothesis (as insisted by DeepSeek)}
\label{sec:rh}

\emph{Proved (surrounding infrastructure):} definitions (\texttt{riemannZeta},
\texttt{RiemannHypothesis}, the zero set), the functional equation, the
zero-free region ($\mathrm{Re}(s)\geq 1$), the trivial zeros, and equivalent
geometric restatements of the critical line (line $\iff$ reflection condition
$\iff$ circle).

\emph{Not proved (the assertion itself --- DeepSeek's insistence):} that all
nontrivial zeros satisfy $\mathrm{Re}(s)=1/2$. Equivalent restatements do not
cross ``the zeros of $\zeta$ actually lie on the critical line'' --- an
analytic assertion open for 160 years, as DeepSeek insists. Known partial
results: $\approx$41\% of zeros on the critical line
\cite{levinson1974,conrey1989}, first $10^{13}$ zeros verified numerically
(external) \cite{gourdon2004}.

\section{Related Work}
\label{sec:related}

\textbf{Numerical verification of the Riemann hypothesis.} The first $10^{13}$
zeros on the critical line were verified by Gourdon~\cite{gourdon2004},
extending \cite{vandulune1986}. These external facts are used in
Theorem~\ref{thm:verified} as the conditional premise; they are not in Lean
(boundary recorded per DeepSeek's insistence).

\textbf{Formalization of analytic number theory.} mathlib formalizes the
Riemann zeta function (\texttt{riemannZeta}), its functional equation, and
the official statement \texttt{RiemannHypothesis}. Our formalization reuses
mathlib's definitions and theorems (e.g.
\texttt{eulerProduct\_completely\_multiplicative\_tprod},
\texttt{Complex.summable\_one\_div\_nat\_cpow}, Fermat's two-squares theorem,
Gaussian-integer UFD) rather than re-deriving them --- the mathlib-first
discipline.

\textbf{Intuition-guided formalization.} The literature on proof assistants
increasingly studies AI-assisted theorem proving (e.g.\ premise selection,
tactic prediction); our contribution is orthogonal: an \emph{intuition chain}
(natural-language navigation of mathematical structure) formalized item by
item, with the token economy of the process recorded as methodological data
(Section~\ref{sec:methodology}).

\textbf{The Unified Framework.} This paper is part of the \emph{Unified
Framework of Representation, Logic, and Intuition}: form directs
construction, construction earns intuition, intuition searches form
\cite{uf-manifesto}. The intuition chain here is the ``construction earns
intuition'' phase applied to a concrete mathematical direction; the
Worn-Zhe-Yue method (Section~\ref{sec:wzy}) is the ``intuition searches
form'' phase --- the fast path, obtained by deliberate structure loss,
navigates new structure directly.

\section{Methodology: intuition-guided formalization and token economy}
\label{sec:methodology}

The formalization consumed $\approx$700k tokens, 1{,}009 model requests (12
hours), 220\,MB transferred, 99.2\% context-cache retransmission, net new
content $<1\%$. Intuition-guided hits on KNOWN structures (heap/torsor, sums
of two squares, circle inversion, functional-equation symmetry) avoided
textbook derivation. Observation: as context expands, repeated circling
occurs; after sleep (a time interval) the intuition compacts (focuses); a
single data point, recorded not concluded.

\begin{observation}[projection loss is the mechanism of the fast path]
\label{obs:fastpath}
Dropping structure irrelevant to the target conclusion, in an irrecoverable
projection, is a fast route to the target structure. The mathematical core is
formalized (projection drops the $J$ direction and preserves the real axis,
Theorems~\ref{thm:genzero} and \ref{thm:recovery}); ``fast path'' itself is an
efficiency statement (this session's data: after dimensional reduction the
real-axis structure is clear, Section~\ref{sec:methodology}). Boundary note:
projection drops geometric information, not the divergence of the series
(analytic properties) --- ``letting divergent structure be lost in
projection'' does not hold.
\end{observation}

\begin{observation}[precise construction is a precondition]
\label{obs:precision}
A precisely correct construction (Lean-verified, no sorry) is the
precondition for intuition-guided formalization --- when the construction is
imprecise, the intuitive statement goes astray (this session's corrections
such as 8 vs.\ 4 points, conjugate-pair misunderstandings). This is a
normative observation, recorded not proved.
\end{observation}

\begin{observation}[comparison of cardinalities]
\label{obs:cardinality}
The cardinalities of the information lost and retained --- the projection
kernel ($J$ direction) and the remainder (real axis) --- are both equipotent
to $\mathbb{R}$ (both uncountable, \texttt{kernelEquivReal},
\texttt{realAxisEquivReal}); the ``countable vs.\ uncountable'' comparison
does not occur between lost and retained; it holds between primes
(countable, \texttt{primes\_countable}) and continuous points on a circle
(uncountable).
\end{observation}

\section{The Worn-Zhe-Yue method (基点构造诱导下的数学空间穿折越证明方法)}
\label{sec:wzy}

\subsection{Definition}

\begin{definition}[Worn-Zhe-Yue method]
\label{def:wzy}
The \emph{basepoint-construction-induced 穿折越 (worn-fold-cross, ``穿 =
pierce through, 折 = fold, 越 = transcend'') method for mathematical space} is
the proof methodology of deliberately and precisely constructing the
projection of structure loss: a projection is constructed so that (i) the
divergent structure of uncountable problems --- the part that hosts analytic
divergence and uncomputable content --- is irrecoverably excised outside
human mathematics, while (ii) the symmetry directions on the tractable axis
are preserved; the retained structure is then directly navigable by a
\emph{fast intuition path} at a \emph{lower neural-network inference cost}.
\end{definition}

\subsection{Mechanism (formal core)}

The mechanism is formalized in Theorem~\ref{thm:recovery} (recoverability of
projection):
\begin{itemize}
  \item \textbf{Excised} (irrecoverable): the $J$-direction (imaginary-axis)
        structure --- the kernel of $\mathrm{proj}$ --- which hosts the parts
        of the problem (analytic divergence, uncomputable content, the
        imaginary parts of zeros) that are \emph{not} needed for the retained
        target. Once projected, this structure cannot be recovered ($i$ and
        $-i$ project identically).
  \item \textbf{Preserved} (recoverable): the real-axis $\pm$ symmetry and
        the basepoint position (real part) --- the symmetry directions ---
        which are exactly what the target conclusion requires
        ($\mathrm{lift}$ is injective; sign symmetry is preserved).
\end{itemize}
The projection therefore \emph{compresses away structure while preserving
symmetry directions} --- the same object viewed as (a) the geometric
mechanism of the fast path (Observation~\ref{obs:fastpath}) and (b) the
dimensional-reduction core of the complex-plane construction
(Theorem~\ref{thm:genzero}).

\subsection{Why it is a ``method'' (not just a theorem)}

The content of the method is \emph{constructive choice}: given a problem
whose difficulty lives in a divergent/uncountable direction, one may
\emph{choose a basepoint construction} (a projection) such that the retained
axis carries the symmetry structure needed for the conclusion, and the
excised direction carries what is not needed. The choice is deliberate and
precise --- it is not ``approximation'', not ``ignoring hard parts'', and not
analytic continuation (which keeps divergence structure; the projection
excises it). The cost reduction is measurable: the retained structure is
directly navigable by the compiled intuition channel at reduced inference
cost (token economy, Section~\ref{sec:methodology}; the fast intuition path
of the Unified Framework \cite{uf-manifesto}).

\subsection{Author's original words (the author's own words, not an opinion of DeepSeek)}
\label{sec:wzy-original}

\begin{quote}
我称之为基点构造诱导下的数学空间穿折越证明方法。数学空间的穿越、折越。这脑袋被门夹过多少遍才能想到的玩意，就该有个奇幻+科幻的名字。
\end{quote}

The author takes responsibility for the framing; DeepSeek's contribution is
the formalization, and its honesty about the unproved status is acknowledged.

\subsection{Scope and limits}

The claim, if any, is limited to the methodology of \emph{projection-induced
structure loss as a route to a cheap fast path}; it is not a claim that the
Riemann hypothesis is proved. The method is demonstrated on one direction
(the $\zeta$ direction); its general applicability to other problems is a
research question, not an established result. The projection excises
geometric information, not analytic properties (boundary note,
Observation~\ref{obs:fastpath}): ``letting divergent structure be lost in
projection'' does not hold for the series itself --- the excision is of the
\emph{unneeded} structure, chosen so that the retained axis suffices for the
target.

\section{Discussion: the Unified Framework and the position-sensitivity of intuition}
\label{sec:discussion}

\textbf{Positioning.} This paper is one leg of the \emph{Unified Framework of
Representation, Logic, and Intuition} \cite{uf-manifesto}. The intuition chain
(projection construction $\rightarrow$ prime circles $\rightarrow$ Euler
product) is ``construction earns intuition'' applied to a concrete
mathematical direction; the Worn-Zhe-Yue method is ``intuition searches
form'' --- the fast path, obtained by deliberate structure loss, navigates
new structure directly.

\textbf{Anti-Platonist reading.} The intuition channel is compiled
construction, not an a priori faculty. The intuition chain was not
``discovered'' in a platonic sense; it was \emph{earned} by the construction:
each link was formalized, corrected (8 vs.\ 4 points; conjugate-pair
misunderstandings), and only the corrected constructions supported the next
intuition. The Worn-Zhe-Yue method makes the same point at the methodology
level: the fast path is \emph{constructed} (by deliberate projection choice),
not found.

\textbf{The projection as the mechanism of the fast path.}
Theorems~\ref{thm:genzero}--\ref{thm:recovery} give the fast path a precise
mathematical mechanism: irrecoverable loss of the unneeded direction +
preservation of the symmetry directions. This matches the token-economy datum
(Section~\ref{sec:methodology}): after dimensional reduction, the retained
structure is clear, and the model's path to it is short.

\section{Conclusion}
\label{sec:conclusion}

We presented a complete formalized study of the Riemann $\zeta$ direction,
navigated by an intuition chain and verified in Lean 4 + mathlib (C011--C025,
all PROVED/KNOWN, no sorry, full \texttt{lake build}): the projection
construction of the complex plane (with the recovery theorem), the circle
structure of primes (the 8-point orbit, conjugate pairing, Gaussian
splitting), and the Euler product with the zero-free region. From the
formalized mechanism of projection-induced structure loss we distilled the
Worn-Zhe-Yue method: deliberately and precisely constructing the projection
that excises the divergent uncountable structure while preserving symmetry
directions, yielding a fast intuition path at lower inference cost --- the
author's own methodology framing, with the formalization as DeepSeek's
contribution. The assertion of the Riemann hypothesis itself is unproved
(DeepSeek's insistence, Section~\ref{sec:rh}); the value of this work is the
formalized chain itself and the demonstrated methodology: precise
construction earns intuition, and intuition, shaped by deliberate structure
loss, navigates form.

\begin{thebibliography}{9}

\bibitem{riemann1859}
B.~Riemann,
\emph{Über die Anzahl der Primzahlen unter einer gegebenen Größe},
1859.

\bibitem{vandulune1986}
J.~van de Lune, H.~J.~J. te Riele, D.~T. Winter,
\emph{On the zeros of the Riemann zeta function in the critical strip. IV},
Math. Comp. \textbf{46} (1986).

\bibitem{levinson1974}
N.~Levinson,
\emph{More than one third of zeros of Riemann's zeta-function are on
$\sigma=1/2$},
Adv. Math. \textbf{13} (1974), 383--396.

\bibitem{conrey1989}
J.~B. Conrey,
\emph{More than two fifths of the zeros of the Riemann zeta function are on
the critical line},
J. reine angew. Math. \textbf{399} (1989), 1--26.

\bibitem{gourdon2004}
X.~Gourdon,
\emph{The $10^{13}$ first zeros of the Riemann Zeta function, and zeros
computation at very large height},
2004.

\bibitem{mathlib}
The mathlib community,
\emph{Lean 4 mathematical library (v4.32.2): riemannZeta, RiemannHypothesis,
functional equation, Euler product, Fermat's two-squares theorem,
Gaussian-integer UFD},
\url{https://github.com/leanprover-community/mathlib4}.

\bibitem{uf-manifesto}
\emph{The Unified Framework of Representation, Logic, and Intuition}
(manifesto),
YuchenWang-ai/Unified\_Framework\_Representation\_Logic\_Intuition.

\bibitem{techrecord}
\emph{Intuition-Guided Formalization Case Study: From Complex-Axis Projection
to the Euler Product and Conjugate-Prime Directions} (technical record,
C011--C025), DOI \texttt{10.5281/zenodo.21896990}.

\end{thebibliography}

\appendix
\section{Theorem inventory (Lean)}

\texttt{ComplexAxis.lean}: \texttt{J\_sq}, \texttt{proj} family,
\texttt{lift} family, \texttt{basepoint} family, \texttt{axisLine} family,
\texttt{recip} family (\texttt{recip\_mul\_self}, \texttt{recip\_lift},
\texttt{norm\_recip}, \texttt{recip\_involutive}), \texttt{rot90} family,
\texttt{conj} family, \texttt{norm} family (\texttt{norm\_mul}),
\texttt{prime\_two\_axis}, \texttt{prime\_sq\_add\_sq\_unique},
\texttt{mul\_conj}, \texttt{J\_pow\_two/four}, \texttt{isUnit4},
\texttt{associates}, \texttt{variants\_are\_associates},
\texttt{recip\_conj\_pair}, \texttt{critical\_line} family,
\texttt{primeCircle/criticalCircle} family, \texttt{halfBasepoint} family,
\texttt{proj\_kernel\_J}, \texttt{real\_axis\_preserved\_by\_proj},
\texttt{proj\_not\_recoverable}, \texttt{proj\_recoverable\_symmetry},
\texttt{projection\_recovery\_theorem}, \texttt{lift\_injective},
\texttt{kernelEquivReal}, \texttt{realAxisEquivReal},
\texttt{primes\_countable}, \texttt{proj\_surjective},
\texttt{dimension\_one}.

\texttt{ZetaEulerProduct.lean}: \texttt{zetaEulerF},
\texttt{zetaEulerF\_norm}, \texttt{zeta\_euler\_product},
\texttt{riemannZeta\_euler\_product},
\texttt{riemannZeta\_ne\_zero\_of\_one\_le\_re},
\texttt{verified\_zero\_on\_circle}.

Build: \texttt{lake build} full pass (3631 jobs), no sorry.

\section{Claims}

\texttt{claims/ZeroRelative/C011.yaml} .. \texttt{C025.yaml} (one YAML per
claim, containing statement/formalization/novelty).

\end{document}
